\Needspace{11\baselineskip}
\section{The matrix structure of the Gauduchon equation}
\label{sec:stw-matrix-structure}

The Hodge star converts the positive $(n-1,n-1)$-form $\Gamma_u$ into a Hermitian matrix whose logarithmic determinant is a concave function of the complex Hessian. Two features will control the largest eigenvalue in the second-order estimate: the positive background gives coercivity, and the torsion coefficients omit exactly those gradient indices that would otherwise carry the largest inverse weight. We derive both properties before differentiating the equation.

For the matrix and derivative calculations, set $v=u/(2\pi)$, so that $\ddc u=\mathrm i\partial\bar\partial v$. All covariant derivatives are taken with respect to the Chern connection $\nabla$ of $\gamma$. In holomorphic coordinates $z^1,\ldots,z^n$, write $v_i=\partial_i v$ and $v_{\bar j}=\partial_{\bar j}v$; indices range from $1$ to $n$, and $\delta_{ij}$ is the Kronecker symbol. Write $\gamma=\mathrm i\gamma_{a\bar b}\,\mathrm{d}z^a\wedge\mathrm{d}\bar z^b$, with repeated indices summed. We denote by $\star_\gamma$ the Hodge star of $\gamma$, and by $\det_\gamma$ the determinant of the Hermitian matrix of a real $(1,1)$-form relative to $\gamma$. Write the torsion of $\nabla$ as $T_{ab}^{c}$ and lower the last index with $\gamma$. Thus
\[
 T_{ab\bar c}=\partial_a\gamma_{b\bar c}-\partial_b\gamma_{a\bar c},
 \qquad T_{ab\bar c}=-T_{ba\bar c}.
\]
Let $Q_0=\star_\gamma\gamma_0^{n-1}/(n-1)!$ and
\[
 Z=\frac1{(n-1)!}\star_\gamma
       \operatorname{Re}\bigl(\mathrm i\partial v\wedge
                    \bar\partial\gamma^{n-2}\bigr).
\]
Here the real $(1,1)$-forms on the right are identified with their Hermitian coefficient matrices relative to $\gamma$. Set
\[
 \chi=(\operatorname{tr}_\gamma Q_0)\gamma-(n-1)Q_0,
 \quad W=(\operatorname{tr}_\gamma Z)\gamma-(n-1)Z,
 \quad g=\chi+\mathrm i\partial\bar\partial v+W.
\]
The linear map $P_\gamma(A)=((\operatorname{tr}_\gamma A)\gamma-A)/(n-1)$ on real $(1,1)$-forms, or Hermitian matrices, is invertible, and the equation of \cref{thm:stw-volume} becomes
\begin{equation}\label{eq:stw-cofactor-matrix-scaling}
 \begin{aligned}
 Q_u&:=P_\gamma(g)=Q_0+P_\gamma(\mathrm i\partial\bar\partial v)+Z>0,\\
 \log\det_\gamma Q_u&=h=(n-1)(F+b).
 \end{aligned}
\end{equation}
Here $Q_u=\star_\gamma\Gamma_u/(n-1)!$ is the coefficient matrix of $\Gamma_u$. The tensor $g$ need not be positive. To check the Hessian term, diagonalize $\mathrm i\partial\bar\partial v$ in a unitary frame: the coefficient of $\star_\gamma(\mathrm i\partial\bar\partial v\wedge\gamma^{n-2})/(n-1)!$ in each diagonal position is the sum of the other diagonal entries divided by $n-1$. If $\eta_u$ has eigenvalues $a_i$ relative to $\gamma$, then $Q_u$ has eigenvalues $\prod_{j\ne i}a_j$, whose product is $(\prod_i a_i)^{n-1}$. This proves the determinant identity.

\Needspace{20\baselineskip}
\begin{lemma}\label{lem:stw-torsion-coefficients}
At a point in a unitary frame,
\begin{align}
 2(n-1)Z_{i\bar j}
 ={}&\delta_{ij}\sum_{p,k}v_p\overline{T_{pk\bar k}}
       -v_i\sum_k\overline{T_{jk\bar k}}
       -\sum_kv_k\overline{T_{kj\bar i}}\notag\\
 &+\delta_{ij}\sum_{p,k}v_{\bar p}T_{pk\bar k}
       -v_{\bar j}\sum_kT_{ik\bar k}
       -\sum_kv_{\bar k}T_{ki\bar j}.
 \label{eq:stw-torsion-six-terms}
\end{align}
\Needspace{9\baselineskip}
In particular $Z_{i\bar j}$ contains no $v_j$ or $v_{\bar i}$. Moreover
\begin{equation}\label{eq:stw-diagonal-torsion}
 2(n-1)Z_{i\bar i}
   =\sum_{p\ne i,\ k\ne i}
        \bigl(v_p\overline{T_{pk\bar k}}+v_{\bar p}T_{pk\bar k}\bigr).
\end{equation}
Covariant differentiation preserves these coefficient identities. For example $\nabla_1 Z_{1\bar1}$ contains no pure Hessian $v_{11}$ or its conjugate, and $\nabla_{\bar1}\nabla_1 Z_{1\bar1}$ contains no $v_{11\bar1}$ or its conjugate, nor a pure $v_{11}$ or its conjugate.
\end{lemma}
\begin{proof}
For $n=2$, the defining expression of $Z$ is zero, and so are the formulas below, because the generalized Kronecker symbol with three indices vanishes in two dimensions. Suppose $n\geq3$ and use coordinates that are unitary at the point. Direct antisymmetrization gives
\[
 \bar\partial\gamma
   =-\frac{\mathrm i}{2}\sum_{a,b,c}
        \overline{T_{bc\bar a}}\,\mathrm{d}z^a\wedge \mathrm{d}\bar z^b\wedge \mathrm{d}\bar z^c.
\]
We use $\bar\partial\gamma^{n-2}=(n-2)\bar\partial\gamma\wedge\gamma^{n-3}$. To write the contraction explicitly, put
\[
 \delta^{paj}_{bci}
   =\det\begin{pmatrix}
     \delta_{pb}&\delta_{pc}&\delta_{pi}\\
     \delta_{ab}&\delta_{ac}&\delta_{ai}\\
     \delta_{jb}&\delta_{jc}&\delta_{ji}
   \end{pmatrix}.
\]
Let $A=\star_\gamma(\mathrm i\partial v\wedge\bar\partial\gamma^{n-2})/(n-1)!$. Expanding the $n-3$ remaining copies of $\gamma$ leaves the three complementary indices and gives
\[
 A_{i\bar j}=\frac1{2(n-1)}
    \sum_{p,a,b,c}v_p\overline{T_{bc\bar a}}\,
                         \delta^{paj}_{bci}.
\]
Indeed, the factor is $(n-2)(n-3)!/[2(n-1)!]$, and the determinant records the signs of the remaining holomorphic and antiholomorphic indices. Expanding its six terms and using $T_{bc\bar a}=-T_{cb\bar a}$, we obtain
\[
 A_{i\bar j}=\frac1{n-1}\left(
   \delta_{ij}\sum_{p,k}v_p\overline{T_{pk\bar k}}
   -v_i\sum_k\overline{T_{jk\bar k}}
   -\sum_kv_k\overline{T_{kj\bar i}}\right).
\]
Since $Z=(A+A^*)/2$, where $A^*$ denotes the Hermitian adjoint matrix, this gives \eqref{eq:stw-torsion-six-terms}.

For $i\ne j$, the coefficient of $v_j$ in this expression vanishes because $T_{jj\bar i}=0$, and the coefficient of $v_{\bar i}$ vanishes in the conjugate expression. For $i=j$, cancelling the terms with $p=i$ and $k=i$ gives \eqref{eq:stw-diagonal-torsion}. These are tensor identities valid in every unitary frame. Since the Chern connection preserves the metric, they can be differentiated in a unitary frame that is parallel at the point. For instance, the first derivative of a summand in \eqref{eq:stw-diagonal-torsion} is
\[
 (\nabla_1v_p)\overline{T_{pk\bar k}}
       +v_p\nabla_1\overline{T_{pk\bar k}},\qquad p,k\ne1,
\]
together with its conjugate counterpart. A further derivative in the direction $\bar1$ still has $p\ne1$. This gives the stated exclusions, including the terms in which a derivative falls on a coefficient.
\end{proof}

The second-order estimate uses these exclusions; an arbitrary first-order perturbation could not be substituted for $Z$.

\Needspace{22\baselineskip}
\begin{lemma}\label{lem:stw-logdet-coercivity}
Diagonalize $g$ at a point and order its eigenvalues as $\lambda_1\geq\cdots\geq\lambda_n$. Put
\[
 \mu_i=\frac{\sum_{j\ne i}\lambda_j}{n-1}>0,
 \quad t_i=\mu_i^{-1},\quad
 f_i=\frac1{n-1}\sum_{j\ne i}t_j,
 \quad \mathcal F=\sum_i f_i=\sum_i t_i.
\]
Then $\mu_1\leq\cdots\leq\mu_n$, $f_1\leq\cdots\leq f_n$, and
\begin{gather}
 \sum_i f_i\lambda_i=n,\qquad
 \frac{\mathcal F}{n(n-1)}\leq f_k\leq\frac{\mathcal F}{n-1}
       \quad(k>1),\qquad t_i\leq(n-1)f_1\quad(i>1),
 \label{eq:stw-weight-comparability}\\
 f_n\geq\mathcal F/n.
 \notag
\end{gather}
Fix $\tau,D>0$ and assume $Q_0\geq\tau\gamma$ and $\det P_\gamma(g)\leq D$. For $\lambda_1$ larger than a constant depending only on $n,\tau,D$,
\begin{equation}\label{eq:stw-direct-coercivity}
 \sum_i f_i(\chi_{i\bar i}-\lambda_i)
                \geq\frac\tau2\mathcal F.
\end{equation}
There is also a uniform positive lower bound for $\mathcal F$ depending only on $n,D$.
\end{lemma}
\begin{proof}
The ordering follows from the formula for $\mu_i$, and then from that for $f_i$. The identity $\sum_if_i\lambda_i=n$, which reflects the relation $\log\det P_\gamma(sg)=\log\det P_\gamma(g)+n\log s$, can be checked directly:
\[
 \sum_i f_i\lambda_i
   =\sum_j\frac{t_j}{n-1}\sum_{i\ne j}\lambda_i
   =\sum_jt_j\mu_j=n.
\]
For $k>1$, the sum defining $f_k$ contains the largest weight $t_1\geq\mathcal F/n$. The upper bound for $f_k$ and the estimate for $t_i$ are immediate, and the largest weight $f_n$ is at least the average $\mathcal F/n$.

The map $P_\gamma$ is self-adjoint for the matrix trace pairing. Consequently
\[
 \sum_i f_i\chi_{i\bar i}
   =\sum_i t_i(Q_0)_{i\bar i}\geq\tau\mathcal F.
\]
Moreover $\sum_i\lambda_i=\sum_i\mu_i>\lambda_1>0$, where the strict inequality follows from $\mu_1>0$. Hence $\mu_n>\lambda_1/n$. Since the other $n-1$ numbers $\mu_i$ are at least $\mu_1$,
\[
 \mu_1^{n-1}\mu_n\leq\prod_i\mu_i\leq D,
 \qquad
 \mathcal F\geq\mu_1^{-1}
             \geq\left(\frac{\lambda_1}{nD}\right)^{1/(n-1)}.
\]
For $\lambda_1$ sufficiently large this is at least $2n/\tau$, so $n\leq\tau\mathcal F/2$. Subtracting $\sum_i f_i\lambda_i=n$ therefore proves \eqref{eq:stw-direct-coercivity}. Finally, the arithmetic--geometric mean inequality applied to the $t_i$ gives $\mathcal F\geq nD^{-1/n}$.
\end{proof}

\Needspace{8\baselineskip}
The concavity of the logarithmic determinant will supply a positive term with an explicit formula. For the matrix function $\Phi(g)=\log\det P_\gamma(g)$, the negative second differential in a complex matrix direction $B$ and its adjoint is
\begin{equation}\label{eq:stw-concavity-square}
 -D^2\Phi_g(B,B^*)
   =\operatorname{tr}\bigl(P_\gamma(g)^{-1}P_\gamma(B)
                   P_\gamma(g)^{-1}P_\gamma(B)^*\bigr)\geq0.
\end{equation}
This follows by differentiating the inverse matrix. In a unitary frame diagonalizing $g$, the off-diagonal terms of this expression include
\[
 \frac{|B_{k\bar1}|^2}{(n-1)^2\mu_k\mu_1}.
\]
When $\lambda_k<\lambda_1$, that coefficient is exactly
\[
 \frac{f_k-f_1}{\lambda_1-\lambda_k}
       =\frac1{(n-1)^2\mu_k\mu_1}.
\]
The expression on the right also makes sense when the two eigenvalues coincide. This positive square and the weights $f_i$ enter the largest-eigenvalue estimate in the next section.
